\section{Sobolev input and absolute continuity}\label{sec:extension}

The signed estimate now gives the endpoint theorem in two steps.
First, approximation passes it to Sobolev inputs and yields a finite
distributional derivative measure for the maximal function. Second,
local mean subtraction shows that this measure gives no mass to a
Lebesgue-null set. The second step is needed to obtain a weak gradient
in $L^1$, and it is the reason the finite-band estimate was proved
for signed inputs.

For real $g\in W^{1,1}(\R^2)$ and $0<a\le b<\infty$, continue to write
\[
 S_{a,b}g(x)=\sup_{a\le t\le b}A_tg(x).
\]
The disk integrals define these functions independently of the
representative of $g$. The positive lower radius makes the supremum
finite at every $x$, since $|A_tg(x)|\le\|g\|_1/(\pi a^2)$.
The elementary analytic facts used below, including the exact
Euclidean variation formula, are proved in
Appendix~\ref{app:analytic-tools}.

\begin{lemma}[Sobolev finite bands]\label{lem:sobolev-band}
For every real $g\in W^{1,1}(\R^2)$ and $0<a\le b<\infty$,
$S_{a,b}g$ is globally Lipschitz and
\begin{equation}\label{eq:band-lipschitz}
 \Lip(S_{a,b}g)\le\frac{\|\nabla g\|_1}{\pi a^2}.
\end{equation}
If $b/a$ is a nonnegative integer power of two, then
\begin{equation}\label{eq:sobolev-band}
 \int_{\R^2}|\nabla S_{a,b}g|\,\dd x\le C\|\nabla g\|_1,
\end{equation}
where $C$ is absolute and independent of the band. This includes
$a=b$, when the envelope is one average.
\end{lemma}

\begin{proof}
The translation estimate of Lemma~\ref{lem:appendix-sobolev-calculus}
gives, for $t\ge a$ and $z\in\R^2$,
\begin{align*}
 |A_tg(x+z)-A_tg(x)|
 &\le\frac1{\pi t^2}\int_{B(x,t)}|g(y+z)-g(y)|\,\dd y\\
 &\le\frac{|z|}{\pi a^2}\|\nabla g\|_1.
\end{align*}
Taking suprema over the same radius interval on both sides proves
\eqref{eq:band-lipschitz}. The same comparison gives, for any
two integrable real inputs,
\begin{equation}\label{eq:band-uniform-approximation}
 \|S_{a,b}g-S_{a,b}\widetilde g\|_\infty
       \le\frac{\|g-\widetilde g\|_1}{\pi a^2}.
\end{equation}
Here the inequality $|\sup f_t-\sup h_t|\le\sup|f_t-h_t|$
is applied to bounded families of real numbers.

Choose real $g_j\in C_c^\infty$ converging to $g$ in $W^{1,1}$.
At the fixed lower radius $a$, \eqref{eq:band-uniform-approximation}
gives uniform convergence of the band envelopes. If $b/a$ is a
power of two greater than one, Proposition~\ref{prop:finite-band}
therefore passes to the limit by the local variation inequality
in Lemma~\ref{lem:appendix-vector-measure}:
\[
 \int|\nabla S_{a,b}g|
 \le\liminf_j\int|\nabla S_{a,b}g_j|
 \le C\lim_j\|\nabla g_j\|_1=C\|\nabla g\|_1.
\]
The limit is Lipschitz by \eqref{eq:band-lipschitz}, so its weak
and almost-everywhere classical gradients agree. When $a=b$,
testing the convolution against smooth functions gives
$\nabla A_ag=A_a(\nabla g)$ distributionally; the integral of
its norm is at most $\|\nabla g\|_1$ by Tonelli's theorem.
\end{proof}

\begin{proof}[Proof of Theorem~\ref{thm:main}]
At this point take $g=|f|$. Lemma~\ref{lem:appendix-sobolev-calculus}
shows that $g\in W^{1,1}(\R^2)$ and
\[
 V_g:=\|\nabla g\|_1\le\|\nabla f\|_1.
\]
The planar Sobolev inequality in
Lemma~\ref{lem:appendix-planar-embedding} gives $g\in L^2(\R^2)$.
The strong $L^2$ maximal inequality \cite[(1.2)]{Kinnunen1997}
then yields $Mf=Mg\in L^2(\R^2)$, hence local integrability.
Define
\[
 G_N=S_{2^{-N},2^N}g,\qquad N=1,2,\ldots.
\]
Because $g\ge0$, these nonnegative functions increase to $Mg$
and are bounded above by it. Dominated convergence on every bounded
set implies $G_N\to Mg$ in $L^1_{\mathrm{loc}}$.

\paragraph{A finite derivative measure.}
Lemma~\ref{lem:sobolev-band} bounds $\int|\nabla G_N|$ by
$CV_g$ uniformly in $N$. In particular, for every
$\varphi\in C_c^1(\R^2;\R^2)$,
\[
 \left|-\int Mg\,\operatorname{div}\varphi\right|
 =\lim_N\left|\int\nabla G_N\cdot\varphi\right|
 \le CV_g\|\varphi\|_\infty.
\]
The vector-measure statement of
Lemma~\ref{lem:appendix-vector-measure} supplies a finite vector
Radon measure $H=D(Mg)$, with
\begin{equation}\label{eq:variation-limit}
 |H|(\R^2)\le CV_g.
\end{equation}
It also gives the local estimate
\begin{equation}\label{eq:local-lsc}
 |H|(O)\le\liminf_{N\to\infty}\int_O|\nabla G_N|\,\dd x
 \qquad(O\subset\R^2\text{ open}).
\end{equation}
These conclusions control total variation. We next show that $H$
is absolutely continuous.

Lower-radius Lipschitz truncations appear in Haj{\l}asz and Mal\'y
\cite[Lemma~8]{HajlaszMaly2010} and Kurka
\cite[Lemma~7.1]{Kurka2015}; local/nonlocal splitting with
subtraction of the input mean appears in Saari
\cite[Section~3.1]{Saari2019}. Here the signed band estimate and
Poincar\'e's inequality make the localization cost depend only
on the nearby input gradient. We give that argument in full.

\paragraph{Small and large radii near a null set.}
Fix a compact Lebesgue-null set $K$. The assertion is immediate
if $K$ is empty, so assume it is nonempty. Fix also a dyadic
$t=2^{-L}$ with $L\ge1$, and an open set
\[
 K\subset O\subset\{x:\dist(x,K)<t\}.
\]
For every $N\ge L$, splitting the band at $t$ gives
\[
 G_N=\max\{S_{2^{-N},t}g,\ S_{t,2^N}g\}.
\]
Both functions are Lipschitz. The maximum rule, including the
equality set, in Lemma~\ref{lem:appendix-sobolev-calculus} gives
\[
 |\nabla G_N|
 \le |\nabla S_{2^{-N},t}g|+|\nabla S_{t,2^N}g|
 \quad\text{a.e.}
\]
For the second term, \eqref{eq:band-lipschitz} bounds its integral
over $O$ by
\begin{equation}\label{eq:large-radius-null-cost}
 \int_O|\nabla S_{t,2^N}g|
       \le \frac{|O|}{\pi t^2}V_g.
\end{equation}

For the first term, partition the plane into half-open grid squares
$Q$ of side $t$, and retain those meeting $O$. There are finitely
many, since $O$ is contained in a bounded neighborhood of $K$.
Choose the cutoff of Lemma~\ref{lem:appendix-square}, equal to one
on $3Q$ and supported in $5Q$, and put
\[
 c_Q=\frac1{|5Q|}\int_{5Q}g,\qquad
 g_Q=\chi_Q(g-c_Q).
\]
The product rule and Poincar\'e's inequality give
\begin{equation}\label{eq:localized-input}
 \|\nabla g_Q\|_1\le C\int_{5Q}|\nabla g|.
\end{equation}
Every disk of radius $s\le t$ centered in $Q$ lies in $3Q$.
Consequently
\[
 A_sg_Q(x)=A_sg(x)-c_Q\quad(x\in Q, 0<s\le t),
 \qquad
 S_{2^{-N},t}g_Q=S_{2^{-N},t}g-c_Q\quad\text{on }Q.
\]
The two envelopes have the same gradient almost everywhere on
$Q$. Since $g_Q$ can have either sign, we use precisely the signed
form of Lemma~\ref{lem:sobolev-band}. The ratio of its endpoints
is $2^{N-L}$, including one when $N=L$. It follows that
\begin{align*}
 \int_O|\nabla S_{2^{-N},t}g|
 &\le\sum_{Q\cap O\ne\varnothing}
                \int_Q|\nabla S_{2^{-N},t}g_Q|\\
 &\le C\sum_{Q\cap O\ne\varnothing}\|\nabla g_Q\|_1
 \le C\int_{\{\dist(x,K)<6t\}}|\nabla g(x)|\,\dd x.
\end{align*}
The last step uses the bounded overlap of the squares $5Q$ and
their containment in the $6t$-neighborhood of $K$, both established
in Lemma~\ref{lem:appendix-square}. Grid boundaries are null and
do not affect these gradient integrals.

Combining this estimate with \eqref{eq:large-radius-null-cost}
and then \eqref{eq:local-lsc} yields
\begin{equation}\label{eq:null-localization}
 |H|(O)\le C\int_{\{\dist(x,K)<6t\}}|\nabla g(x)|\,\dd x
                  +C t^{-2}|O|V_g.
\end{equation}
The constants are independent of $N$, $O$, and $t$.

\paragraph{Removing the singular part.}
For the fixed $t$, compactness and Lebesgue nullity of $K$ allow
the open neighborhood $O$ above to have arbitrarily small area.
Since $|H|(K)\le|H|(O)$, \eqref{eq:null-localization} gives
\[
 |H|(K)\le C\int_{\{\dist(x,K)<6t\}}|\nabla g(x)|\,\dd x.
\]
Only now let $t$ decrease to zero through dyadic values. The
neighborhoods decrease to the compact null set $K$, and
$|\nabla g|$ is integrable; the right side tends to zero.
Thus $|H|(K)=0$ for every compact Lebesgue-null set. Inner
regularity of the finite Radon measure $|H|$ implies that it
vanishes on every Borel null set, and hence on every Lebesgue-null
set by containment in a Borel null set.

The Radon--Nikodym conclusion of
Lemma~\ref{lem:appendix-vector-measure} now gives
$H=G\,\dd x$ with $G\in L^1(\R^2;\R^2)$ and
\[
 \int|G|=|H|(\R^2)\le CV_g\le C\int|\nabla f|.
\]
For every smooth compactly supported vector field $\varphi$,
the defining identity of $H$ reads
$-\int Mf\,\operatorname{div}\varphi=\int G\cdot\varphi$.
Thus $G$ is the weak gradient of $Mf$. Since $Mf$ is locally
integrable, this proves $Mf\in W^{1,1}_{\mathrm{loc}}$ and
\eqref{eq:main}, with a globally integrable Euclidean gradient.
\end{proof}

\begin{remark}
The direct absolute-continuity argument uses the signed band
estimate, approximation, and square Poincar\'e localization.
After \eqref{eq:variation-limit}, one could alternatively apply
Lahti and Weigt's centered regularity theorem
\cite[Theorem~1.2]{LahtiWeigt2025} to $g=|f|$ on $\R^2$.
The input is globally integrable and has finite variation, and
$Mg\in BV_{\mathrm{loc}}$ has already been proved. Its input
jump and Cantor derivative measures vanish, so the theorem's
additional condition on the Cantor part is vacuous; its conclusion
gives local Sobolev regularity. This optional route is used only
after the variation bound and supplies no part of the finite-band
estimate.
\end{remark}

\begin{corollary}[Bounded-variation inputs]\label{cor:bv-input}
Let $BV(\R^2)$ denote the globally integrable functions whose
distributional derivatives have finite total variation. For every
real $f\in BV(\R^2)$,
\[
 Mf\in L^2(\R^2)\cap BV_{\mathrm{loc}}(\R^2),
 \qquad |D(Mf)|(\R^2)\le C|Df|(\R^2),
\]
with the constant from Theorem~\ref{thm:main}.
Writing $P(F)=|D\1_F|(\R^2)$ for perimeter, every measurable
set $E$ with $|E|<\infty$ and $P(E)<\infty$ satisfies
\[
 \int_0^1 P(\{M\1_E>t\})\,\dd t\le CP(E).
\]
These superlevel sets have finite measure for every $t>0$ and
finite perimeter for almost every $t>0$.
\end{corollary}

\begin{proof}
Let $f_\varepsilon=f*\rho_\varepsilon$ for a smooth nonnegative
compactly supported probability approximate identity. Then
$f_\varepsilon\in W^{1,1}(\R^2)$, $f_\varepsilon\to f$ in $L^1$,
and
\[
 \nabla f_\varepsilon=(Df)*\rho_\varepsilon,\qquad
 \|\nabla f_\varepsilon\|_1\le |Df|(\R^2).
\]
The variation formula and lower semicontinuity show that these
gradient norms tend to $|Df|(\R^2)$, so the approximation is strict
in $BV$. Lemma~\ref{lem:appendix-planar-embedding} bounds
$\|f_\varepsilon\|_2$ uniformly; an almost-everywhere convergent
subsequence and Fatou's lemma give $f\in L^2$. The convolution
approximate identity therefore also converges to $f$ in $L^2$.
The pointwise inequality
$|Mf_\varepsilon-Mf|\le M(f_\varepsilon-f)$ and the strong $L^2$
maximal inequality yield $Mf\in L^2$ and
$Mf_\varepsilon\to Mf$ in $L^2$, hence locally in $L^1$.
Theorem~\ref{thm:main} bounds the global weak-gradient norm of
$Mf_\varepsilon$ by $C|Df|(\R^2)$.
Lemma~\ref{lem:appendix-vector-measure} now gives the asserted
finite derivative measure and local $BV$ regularity.

For $u=M\1_E$ we have $0\le u\le1$. Apply the $BV$ coarea formula
\cite{EvansGariepy2015} on $B(0,R)$ and then use monotone exhaustion
as $R\to\infty$ to obtain
\[
 \int_0^1 P(\{u>t\})\,\dd t
 =|Du|(\R^2)\le CP(E).
\]
Also $|\{u>t\}|\le t^{-2}\|u\|_2^2<\infty$ for every $t>0$.
\end{proof}

This conclusion does not assert global integrability of $Mf$ or
absolute continuity of its derivative measure for arbitrary $BV$ input.
